% SHARD-ID: AISM-28-HCB4-CANONICAL-INVERSE
% SHARD-TITLE: The canonical Ha inverse estimate
% SHARD-SUMMARY: Reproduces lem-hcb4-canonical-inverse, the af-validated complete bijectivity of the two special Ha maps with amplified inverses within O(e) of the canonical ones.
% SHARD-SUMMARY: Explains the canonical inverse bound, the quantitative Neumann step, and the closed-corner completeness branch required for the Banach target.
% SHARD-KEYWORDS: lem-hcb4-canonical-inverse, af validated, Neumann series, complete bijectivity, Banach target, closed corner
\section{The canonical Ha inverse estimate}\label{sec:hcb4-canonical-inverse}

\begin{lemma}[\texttt{lem-hcb4-canonical-inverse}]\label{lem:hcb4-canonical-inverse}
There are universal $C_{\mathrm{sp,inv}}<\infty$ and $\comb_{\mathrm{sp,inv}}>0$ such that
every H-CB datum with $\comb\le \comb_{\mathrm{sp,inv}}$ has the special maps
$\Ha{Q}{P}{Q}$ and $\Ha{Q}{Q}{P}$ completely bijective, with their amplified inverses
differing from the corresponding canonical inverses by at most
\begin{equation}\label{eq:ci-main}
  \bigl\lVert\bigl(\Han{Q}{P}{Q}{n}\bigr)^{-1}-J_{P,Q,n}^{-1}\bigr\rVert
  \;\le\;C_{\mathrm{sp,inv}}\,\comb ,
\end{equation}
and likewise for the row pair $\bigl(\Han{Q}{Q}{P}{n},J_{Q,P,n}\bigr)$.
\emph{Transcription note.} ``Completely bijective'' means bijective at \emph{every}
amplification level $n\ge1$, not merely at level one; and the estimate is uniform in $n$.
Unlike the diagonal and off-diagonal inverse statements of
\S\S\ref{sec:hcb3-diagonal-inverse}--\ref{sec:hcb3-offdiagonal-inverse}, this lemma is
\emph{unconditional}: no level-one hypothesis is assumed, because the comparison map
$J_n$ is bijective by construction and the smallness of $\comb$ does the rest.
\end{lemma}

\noindent\textbf{Registry contract (verbatim).}
\contractquote{Canonical Ha inverse estimate: there are universal C_sp,inv < infinity and e_sp,inv > 0 such that every H-CB datum with e <= e_sp,inv has the special maps Ha^Q_{P,Q} and Ha^Q_{Q,P} completely bijective, with their amplified inverses differing from the corresponding canonical inverses by at most C_sp,inv*e.}

\paragraph{Proof (prose account of the af-validated tree).}
The mathematics is a perturbation of an invertible map by a small one; the length of the
tree is due almost entirely to justifying that the target space is Banach, which is where
the Neumann series lives.

\emph{The ledger and the canonical inverse.} Let $C_J,\comb_J$ come from
\texttt{lem-hcb4-canonical-gram} (\S\ref{sec:hcb4-canonical-gram}) and
$C_{\mathrm{sp}},\comb_{\mathrm{sp}}$ from \texttt{lem-hcb4-canonical-closeness}
(\S\ref{sec:hcb4-canonical-closeness}), and put
\begin{equation}\label{eq:ci-ledger}
  \comb_{\mathrm{sp,inv}}:=\min\Bigl\{\comb_J,\ \comb_{\mathrm{sp}},\
    \tfrac1{4(C_J+C_{\mathrm{sp}}+1)}\Bigr\},
  \qquad
  C_{\mathrm{sp,inv}}:=4\bigl(C_{\mathrm{sp}}+1\bigr).
\end{equation}
The last entry of the threshold forces both $C_J\comb\le\tfrac14$ and
$C_{\mathrm{sp}}\comb\le\tfrac14$ simultaneously. For either special ordered corner and
every $n$, the canonical map $J_n$ is the amplified coefficient-space identification of
\texttt{def-canonical-corner-identifications} --- the column map, or the row map defined as
its adjoint --- hence an algebraic linear bijection onto the indicated operator corner,
with the corresponding canonical inverse. Its inverse is bounded: the lower half of the
Gram estimate reads $\lVert J_nZ\rVert\ge(1-C_J\comb)\nrm{Z}$, and applying it to
$Z=J_n^{-1}Y$ with $C_J\comb\le\tfrac14$ gives
\begin{equation}\label{eq:ci-jinv}
  \bigl\lVert J_n^{-1}\bigr\rVert\;\le\;\tfrac43 .
\end{equation}

\emph{The quantitative Neumann step.} The abstract step used twice below is: if
$J:X\to Y$ is a bounded bijection onto a \emph{Banach} space $Y$ with
$\lVert J^{-1}\rVert\le\tfrac43$, and $A:X\to Y$ satisfies
$\lVert A-J\rVert\le C_{\mathrm{sp}}\comb$ with $C_{\mathrm{sp}}\comb\le\tfrac14$, then
$A$ is bijective and $\lVert A^{-1}-J^{-1}\rVert\le\tfrac83C_{\mathrm{sp}}\comb$. Put
$T:=(A-J)J^{-1}:Y\to Y$, so that $A=(I_Y+T)J$ and
$\lVert T\rVert\le\lVert A-J\rVert\lVert J^{-1}\rVert\le\tfrac43C_{\mathrm{sp}}\comb
\le\tfrac13$. Since $Y$ is Banach so is $B(Y)$, and the partial sums
$S_N=\sum_{k=0}^{N}(-T)^{k}$ are Cauchy there by
$\sum_{k\ge0}\lVert T\rVert^{k}\le(1-\lVert T\rVert)^{-1}$; passing to the limit in the
finite geometric identities $(I_Y+T)S_N=S_N(I_Y+T)=I_Y-(-T)^{N+1}$ and using
$\lVert T\rVert^{N+1}\to0$ identifies the limit as the two-sided inverse, with
$\lVert(I_Y+T)^{-1}\rVert\le\tfrac32$. Then $A^{-1}=J^{-1}(I_Y+T)^{-1}$ and
$(I_Y+T)^{-1}-I_Y=-(I_Y+T)^{-1}T$ give
\begin{equation}\label{eq:ci-neumann}
  \bigl\lVert A^{-1}-J^{-1}\bigr\rVert
  \;\le\;\tfrac43\cdot\tfrac32\cdot\tfrac43\,C_{\mathrm{sp}}\comb
  \;=\;\tfrac83\,C_{\mathrm{sp}}\,\comb .
\end{equation}

\emph{Why the targets are Banach.} This is not a formality, and the tree does not treat it
as one. In an H-CB datum the ambient $\mathcal A$ is complete by
\texttt{def-extended-epsilon-cstar-algebra}, and \texttt{def-compressed-corner} states that
$\Co{P}{Q}$ is an exact bounded idempotent with
$\Cnr{P}{Q}=\Img\Co{P}{Q}=\Ker(1-\Co{P}{Q})$; a kernel of a continuous map is closed, and a
closed subspace of a Banach space is Banach, so $\Cnr{P}{Q}$ is complete in the
\emph{inherited} norm. The norm actually carried by the target is the Euclidean one, and
the two are equivalent by the Gram estimate at $n=1$: there $J_{P,Q,1}(Z)c=Zc$, so
$\lVert J_{P,Q,1}(Z)\rVert=\lVert Z\rVert_{\mathrm{Euc}}$ and
$(1-C_J\comb)\nrm{Z}\le\lVert Z\rVert_{\mathrm{Euc}}\le(1+C_J\comb)\nrm{Z}$ with
$C_J\comb\le\tfrac14$. Hence $(\Cnr{P}{Q},\lVert\cdot\rVert_{\mathrm{Euc}})$ is complete,
its finite Hilbertian sum $\mathbb C^{n}\otimes\Cnr{P}{Q}$ is complete, and since
$B(E,F)$ is Banach whenever $F$ is, both actual targets
$B(\mathbb C^{n},\mathbb C^{n}\otimes\Cnr{P}{Q})$ and
$B(\mathbb C^{n}\otimes\Cnr{P}{Q},\mathbb C^{n})$ --- the column and row targets of
\texttt{def-ha-map} and \texttt{def-canonical-corner-identifications} --- are Banach.

\emph{Assembly.} Fix $n\ge1$ and either special corner, and let $A_n$ be
$\Han{Q}{P}{Q}{n}$ or $\Han{Q}{Q}{P}{n}$. Canonical closeness gives
$\lVert A_n-J_n\rVert\le C_{\mathrm{sp}}\comb$, \eqref{eq:ci-jinv} gives the bijection $J_n$
with $\lVert J_n^{-1}\rVert\le\tfrac43$, and \eqref{eq:ci-ledger} gives
$C_{\mathrm{sp}}\comb\le\tfrac14$; so the Neumann step applies and
$\lVert A_n^{-1}-J_n^{-1}\rVert\le\tfrac83C_{\mathrm{sp}}\comb\le
4(C_{\mathrm{sp}}+1)\comb=C_{\mathrm{sp,inv}}\comb$. The constants do not depend on $n$ or
on the corner, so both special maps are completely bijective with the uniform estimate
\eqref{eq:ci-main}.

\medskip
\noindent\textbf{Status.} \texttt{proved}; \texttt{af: validated} (T0). Workspace
\texttt{proofs/lem-hcb4-canonical-inverse}: 19 nodes, all \texttt{validated}, root
\texttt{validated}, taint clean. Six of the nineteen exported nodes form the completeness
branch.

\noindent\textbf{Role in the chain.} Registry imports:
\texttt{lem-hcb4-canonical-gram},
\texttt{lem-hcb4-canonical-closeness},
\texttt{lem-compcb-corner-algebra} and
\texttt{lem-hcb3-uniform-square-lower}; the exported tree visibly consumes the first two.
Its \texttt{defs} line records \texttt{def-ha-map}, \texttt{def-hcb-datum},
\texttt{def-canonical-corner-identifications}, \texttt{def-compressed-corner} (the
range/kernel identity) and \texttt{def-extended-epsilon-cstar-algebra} (completeness of
$\mathcal A$); the last two are consumed essentially in the validated
closed-corner completeness branch. Downstream, this lemma is one of
the sixteen declared dependencies of \texttt{conj-hcb} (\S\ref{sec:hcb}); the parent's own
exported tree cites eight suppliers directly and this is not among them, so the parent
consumes it through the registry rather than in its validated derivation. Nothing here
bears on \texttt{op-classical}, which is now discharged at the af-validated rung via Route F (\S\ref{sec:status-outlook}); sharpness remains with \texttt{ex-hume}.
